\documentclass[11pt]{article}
% === core_packages.tex ===
\usepackage{amsmath, amssymb, amsthm, amsfonts}
\usepackage{geometry}
\usepackage{hyperref}
\usepackage{graphicx}
\usepackage{physics}
\usepackage{booktabs}
\usepackage{listings}
\usepackage{bookmark}
\usepackage{amscd}
\usepackage{tikz-cd}
\usepackage{mathtools}
\usepackage{xspace}
\usepackage[backend=biber, style=numeric]{biblatex}
\usepackage{tcolorbox}
\tcbuselibrary{theorems}
\usepackage{parskip}
\usepackage{tikz}
\usepackage{etoolbox}
\usetikzlibrary{arrows.meta, shapes, positioning, calc, cd, shapes.geometric, backgrounds}
\usepackage[en-US]{datetime2}
\usepackage{array}


% --- Math & Ontology Macros ---
\newcommand{\entropy}{S}
\newcommand{\opponent}{opponent}
\newcommand{\aside}[1]{[\emph{\opponent}: ``#1'']}
\newcommand{\paper}[1]{[\emph{#1}: ``#1'']}
\newcommand{\iame}{Real\textsuperscript{e}}
\newcommand{\iameverything}{Real\textsuperscript{everything}}
\newcommand{\fabric}{voxel}
\newcommand{\Fabric}{Voxel}
\newcommand{\pdflink}[1]{}



% --- Global Layout Defaults ---
\geometry{letterpaper, margin=1in}
\hypersetup{
    colorlinks=true,
    linkcolor=black,
    citecolor=black,
    filecolor=magenta,
    urlcolor=blue
}
\hfuzz=5pt \vfuzz=5pt \hbadness=10000 \vbadness=10000

\author{\iame}


\usepackage{tocloft}

% --- Article Specific Theorems ---
\newtheorem{lemma}{Lemma}
\theoremstyle{definition}
\newtheorem{axiom}{Axiom}
\newtheorem{theorem}{Theorem}
\newtheorem{corollary}{Corollary}
\newtheorem{definition}{Definition}
\newtheorem{proposition}{Proposition}
\newtheorem{conjecture}{Conjecture}


\makeatletter
\let\remark\@undefined
\let\endremark\@undefined
\makeatother
\theoremstyle{remark}
\newtheorem*{remark}{Remark}

\newtcbtheorem[number within=section]{principle}{Principle}%
{colback=blue!5,colframe=blue!50!black,fonttitle=\bfseries}{principle}

% --- Article TOC Layout ---
\setlength{\cftbeforesecskip}{1.0em}
\renewcommand{\cftsecleader}{\cftdotfill{\cftdotsep}}
\renewcommand{\cftsubsecleader}{\cftdotfill{\cftdotsep}}


\addbibresource{../references.bib}

\title{The Universal Boson Theorem:\\
Particle Species from Codec Geometry}
\author{Juha Meskanen}
\date{June 2026}

\begin{document}

\maketitle

\begin{abstract}
Paper~VII established that a photon-like boson emerges as the compression
residual of a two-site fermion hop, with amplitude
$\|B(\theta)\| = \sin(2\theta)/\sqrt{2}$.
The present paper extends this result to $n$-site fermion systems and
proves an exact universal theorem:
for any fermion in equal superposition across $n$ sites with arbitrary
phases, the Frobenius norm of the boson matrix satisfies
\[
  \|B\| = \sqrt{\frac{n-1}{n}}.
\]
This result is phase-independent and holds for all $n \geq 2$.
The phase \emph{structure} of the wavefunction --- specifically the
winding number $m$ of the phase pattern around the $n$-site ring ---
determines the \emph{spin} of the emergent boson while leaving its norm
unchanged.
Together, the pair $(n, m)$ classifies the entire known boson spectrum:
the photon $(n{=}2, m{=}1)$, the gluon $(n{=}3, m{=}1)$, a scalar
Higgs candidate $(n{=}3, m{=}0)$, and a spin-2 graviton candidate
$(n{=}4, m{=}2)$.
Fractional electric charge emerges exactly as $1/n$ per colour site,
without postulate.
A codec argument for colour confinement is given based on the mixed
symmetry of the three-site boson matrix.
\end{abstract}

\tableofcontents
\newpage

%------------------------------------------------------------
\section{Introduction}
\label{sec:intro}

The informational framework of Papers~I--VII
\cite{meskanen2001,meskanen2026iv,meskanen2026v,meskanen2026vi,meskanen2026vii}
treats the universe as a finite static bitstring of $n \approx 184$ bits.
Internal observers traverse the $2^n$ configurations along paths of
minimum description length.
The quantum wavefunction is the compression codec; bosons are its
off-diagonal residuals.

Paper~VII showed that a two-site fermion hop produces a boson matrix $B$
with $\|B\| = 1/\sqrt{2}$, purely antisymmetric structure, and a universal
$-\pi/2$ phase --- the defining properties of a virtual photon.
That result was derived analytically and confirmed numerically to
floating-point precision.

The natural question is: what happens for three or more sites?
Quarks are not two-state objects.
They carry colour charge, they come in triplets, and their mediating
bosons (gluons) have a richer structure than photons.
If the codec framework is to reproduce the Standard Model, it must
generate this richer structure from first principles.

The present paper answers this question with a single theorem whose
proof occupies three lines, followed by a classification of the
resulting particle spectrum.

%------------------------------------------------------------
\section{Setup}
\label{sec:setup}

We follow the conventions of Paper~VII exactly.
A \emph{fermion} is a localised excitation occupying one site of an
$n$-site system.
Its state is encoded as a complex wavefunction
$\psi \in \mathbb{C}^n$, normalised to $\|\psi\|^2 = 1$.
The \emph{density matrix} is $\rho = |\psi\rangle\langle\psi|$, with
components $\rho_{ij} = \psi_i \psi_j^*$.
The \emph{boson matrix} is
\[
  B = \rho - \operatorname{diag}(\rho),
\]
the off-diagonal part of $\rho$.
Its Frobenius norm $\|B\| = \sqrt{\sum_{i \neq j} |B_{ij}|^2}$ measures
the total boson amplitude.

An \emph{$n$-site equal superposition} is a state of the form
\[
  \psi_k = \frac{1}{\sqrt{n}}\,e^{i\phi_k}, \qquad k = 0, \ldots, n-1,
\]
for arbitrary real phases $\phi_k$.
Each site carries probability $|\psi_k|^2 = 1/n$.

%------------------------------------------------------------
\section{The Universal Boson Theorem}
\label{sec:theorem}

\begin{theorem}[Universal Boson Norm]
\label{thm:main}
Let $\psi$ be any $n$-site equal superposition with arbitrary phases
$\phi_0, \ldots, \phi_{n-1}$.
Then
\[
  \|B\| = \sqrt{\frac{n-1}{n}}.
\]
\end{theorem}

\begin{proof}
The density matrix entries are
\[
  \rho_{ij} = \psi_i\psi_j^* = \frac{1}{n}\,e^{i(\phi_i - \phi_j)}.
\]
The diagonal entries are $\rho_{ii} = 1/n$.
The boson matrix has entries $B_{ij} = \rho_{ij}$ for $i \neq j$ and
$B_{ii} = 0$.
Therefore
\[
  \|B\|^2
  = \sum_{i \neq j} |B_{ij}|^2
  = \sum_{i \neq j} \frac{1}{n^2}\,\bigl|e^{i(\phi_i-\phi_j)}\bigr|^2
  = \sum_{i \neq j} \frac{1}{n^2}
  = n(n-1)\cdot\frac{1}{n^2}
  = \frac{n-1}{n}.
\]
Taking the square root gives the result. \qed
\end{proof}

\begin{remark}
The proof uses only the normalisation $|\psi_k| = 1/\sqrt{n}$ and the
definition $B = \rho - \operatorname{diag}(\rho)$.
It is independent of all phases $\phi_k$, of the system size $n$, of any
spatial arrangement of sites, and of any overall phase factor.
\end{remark}

\begin{corollary}[Recovery of Paper VII]
For $n = 2$ with $\psi = (1/\sqrt{2})(e_0 + ie_1)$,
Theorem~\ref{thm:main} gives $\|B\| = \sqrt{1/2} = 1/\sqrt{2}$,
recovering the exact result of Paper~VII.
\end{corollary}

\begin{corollary}[Asymptotic behaviour]
As $n \to \infty$, $\|B\| \to 1$.
A fermion delocalised over infinitely many sites produces a boson of unit
amplitude.
The vacuum limit $n = 1$ gives $\|B\| = 0$: a localised fermion produces
no boson, consistent with Pauli exclusion (Paper~VII, Result~1).
\end{corollary}

Table~\ref{tab:norms} lists $\|B\|$ for small $n$.

\begin{table}[h]
\centering
\begin{tabular}{cccc}
\toprule
$n$ & $\|B\| = \sqrt{(n-1)/n}$ & Decimal & Particle candidate \\
\midrule
1 & $0$           & $0$      & (localised, no boson) \\
2 & $1/\sqrt{2}$  & $0.7071$ & Photon \\
3 & $\sqrt{2/3}$  & $0.8165$ & Gluon / Higgs \\
4 & $\sqrt{3/4}$  & $0.8660$ & Graviton candidate \\
5 & $\sqrt{4/5}$  & $0.8944$ & --- \\
\bottomrule
\end{tabular}
\caption{Universal boson norms for $n$-site equal superpositions.}
\label{tab:norms}
\end{table}

%------------------------------------------------------------
\section{Spin from Winding Number}
\label{sec:spin}

Theorem~\ref{thm:main} shows that $\|B\|$ depends only on $n$.
The \emph{identity} of the boson --- its spin, symmetry, and propagator
structure --- is determined by the phase pattern $\{\phi_k\}$.

The natural phase patterns on an $n$-site ring are the discrete Fourier
modes
\[
  \phi_k^{(m)} = \frac{2\pi m k}{n}, \qquad m = 0, 1, \ldots, \lfloor n/2 \rfloor,
\]
where $m$ is the \emph{winding number}: the number of complete phase
revolutions as $k$ runs from $0$ to $n-1$.
These are the minimum-spectral-complexity states on the ring, ranked by
$m$ in ascending order of complexity cost $C_s$.

\begin{definition}[Codec spin]
The \emph{codec spin} of an $n$-site boson with winding number $m$ is $m$.
\end{definition}

Numerical verification (Table~\ref{tab:spin}) confirms the symmetry
classification:

\begin{table}[h]
\centering
\begin{tabular}{ccclll}
\toprule
$n$ & $m$ & $\|B\|$ & Symmetry of $B$ & Codec spin & Particle \\
\midrule
2 & 0 & $1/\sqrt{2}$ & symmetric     & 0 & scalar   \\
2 & 1 & $1/\sqrt{2}$ & antisymmetric & 1 & \textbf{Photon} \\
3 & 0 & $\sqrt{2/3}$ & symmetric     & 0 & \textbf{Higgs candidate} \\
3 & 1 & $\sqrt{2/3}$ & mixed         & 1 & \textbf{Gluon} \\
4 & 0 & $\sqrt{3/4}$ & symmetric     & 0 & scalar   \\
4 & 1 & $\sqrt{3/4}$ & mixed         & 1 & ---      \\
4 & 2 & $\sqrt{3/4}$ & symmetric     & 2 & \textbf{Graviton candidate} \\
\bottomrule
\end{tabular}
\caption{Spin classification of $n$-site bosons by winding number $m$.
Symmetry refers to $B$ vs $B^\top$: antisymmetric means $B = -B^\top$
(spin-1 propagator); symmetric means $B = B^\top$ (spin-0 or spin-2);
mixed means neither.}
\label{tab:spin}
\end{table}

%------------------------------------------------------------
\section{The Three-Site Boson: Quarks and Gluons}
\label{sec:threesite}

\subsection{Fractional charge}

A three-site equal superposition assigns probability $|\psi_k|^2 = 1/3$
to each site.
If electric charge is identified with the per-site probability weight,
a three-site fermion carries charge $1/3$ per colour component.
This is the electric charge of the down quark ($-1/3$) and, with two
windings, the up quark ($+2/3$).
Fractional charge is not postulated; it is the exact probability weight
of a three-site codec state.

\subsection{The gluon as compression residual}

For the winding-1 state
\[
  \psi = \frac{1}{\sqrt{3}}\bigl(e_0 + i\,e_1 - e_2\bigr),
\]
the boson matrix has entries
\[
  B_{01} = -\frac{i}{3}, \quad
  B_{02} = -\frac{1}{3}, \quad
  B_{12} = -\frac{i}{3},
\]
and their conjugate transposes, with $|B_{ij}| = 1/3$ for all $i \neq j$.
The three off-diagonal magnitudes are exactly equal.
This is \emph{colour symmetry}: the boson couples equally to all three
colour pairs.
No $\mathrm{SU}(3)$ symmetry was imposed; it emerged from the equal
superposition requirement combined with Theorem~\ref{thm:main}.

\subsection{Colour confinement from mixed symmetry}

The photon boson matrix ($n=2$, $m=1$) is purely antisymmetric:
$B = -B^\top$.
Antisymmetry is preserved under spatial inversion, allowing the boson
to propagate freely to infinity and produce a long-range force.

The gluon boson matrix ($n=3$, $m=1$) is neither symmetric nor
antisymmetric.
Mixed-symmetry objects cannot propagate as free asymptotic states under
a codec that selects for minimum description length: a mixed-symmetry
residual always requires reference to the fermion configuration from
which it arose, because its symmetry class is not self-contained.
In physical language, the gluon cannot exist as an isolated particle;
it is confined to the colour-charged system that generated it.
Colour confinement is a symmetry constraint of the codec, not a
dynamical mechanism to be derived from a potential.

%------------------------------------------------------------
\section{The Four-Site Boson: Graviton Candidate}
\label{sec:graviton}

For $n=4$, $m=2$ (two full phase windings):
\[
  \psi = \frac{1}{2}(e_0 + i\,e_1 - e_2 - i\,e_3).
\]
The boson matrix $B$ is symmetric ($B = B^\top$) and has codec spin 2.
This is the only boson in the classification with spin 2, consistent
with the graviton requirement from general relativity.

Moreover, the four-site fermion lies at a natural boundary in the codec
hierarchy.
The two-site fermion generates the $\mathrm{U}(1)$ gauge boson (photon).
The three-site fermion generates the colour-triplet gauge boson (gluon).
The four-site fermion generates the spin-2 metric perturbation (graviton).
The pattern $n = 2, 3, 4$ mirrors the gauge group rank sequence
$\mathrm{U}(1), \mathrm{SU}(3), \text{gravity}$, with $n-1$ giving
the number of independent colour degrees of freedom.

%------------------------------------------------------------
\section{The Particle Spectrum in Log-Mass Space}
\label{sec:spectrum}

The spectral complexity $C_s$ of Paper~VI assigns an information cost to
each fermionic state proportional to its dominant wavenumber $\omega$.
Under Solomonoff induction, $P(\psi) \propto 2^{-C_s(\psi)}$, so the
observable particle spectrum is ordered by ascending $C_s$.

Computing $\log_2(m/m_e)$ for the known massive particles reveals a
near-integer sequence (Table~\ref{tab:masses}), where the step size of
approximately 4 bits corresponds to the minimum fermionic codec unit:
one fermion hop requires four binary choices (which site, which frame,
real or imaginary component, sign of phase), giving $2^4 = 16$
configurations.

\begin{table}[h]
\centering
\begin{tabular}{lccc}
\toprule
Particle & Mass (MeV) & $\log_2(m/m_e)$ & Nearest integer \\
\midrule
Electron & $0.511$    & $0.000$ & $0$  \\
Muon     & $105.66$   & $7.692$ & $8$  \\
Tau      & $1776.86$  & $11.764$& $12$ \\
W boson  & $80{,}400$ & $17.264$& $17$ \\
Z boson  & $91{,}200$ & $17.445$& $17$ \\
Higgs    & $125{,}100$& $17.901$& $18$ \\
\bottomrule
\end{tabular}
\caption{Known particle masses in log$_2$ space relative to the electron.
The near-integer values suggest a discrete complexity ladder with step
size $\approx 4$ bits.
Deviations from exact integers are interpreted as compression gains:
the actual particle is slightly cheaper to describe than the
nearest-integer prediction, consistent with Solomonoff induction
favouring more compressible states.}
\label{tab:masses}
\end{table}

The lepton generations (electron, muon, tau) are identified as the same
fermionic knot topology at three successive compression levels, separated
by approximately 4 bits per generation.
The massive electroweak bosons (W, Z, Higgs) appear at bits 17--18,
consistent with their identification as cross-peak lognormal residuals:
compression residuals from fermion hops that cross the maximum of the
microstructure probability distribution.

%------------------------------------------------------------
\section{Summary of Results}
\label{sec:results}

Four results are established, of which the first is analytically exact
and the remainder follow from it by classification:

\begin{enumerate}
\item \textbf{Universal Boson Theorem} (Theorem~\ref{thm:main}).
  For any $n$-site equal superposition with arbitrary phases,
  $\|B\| = \sqrt{(n-1)/n}$.
  This is exact, phase-independent, and encoding-independent.

\item \textbf{Spin from winding number.}
  The codec spin of the emergent boson equals the winding number $m$
  of the phase pattern around the $n$-site ring.
  The pair $(n, m)$ uniquely classifies each boson species.

\item \textbf{Fractional charge from site count.}
  A fermion in equal superposition across $n$ sites carries charge $1/n$
  per colour component.
  For $n = 3$: charge $1/3$, identifying the state as a quark.

\item \textbf{Colour confinement from mixed symmetry.}
  The three-site ($n=3$, $m=1$) boson matrix has mixed symmetry, which
  prevents it from existing as a free asymptotic state.
  Confinement is a structural property of the codec, not a dynamical one.
\end{enumerate}

%------------------------------------------------------------
\section{Open Problems}
\label{sec:open}

\begin{enumerate}
\item \textbf{Why three generations?}
  The lepton ladder has three rungs (0, 8, 12 bits).
  The lognormal probability distribution on the rising slope has finite
  support; a fourth generation at $\approx 16$ bits would require a
  fermion complexity cost that places it beyond the observer's existence
  saddle point.
  This argument should be made quantitative.

\item \textbf{W and Z from first principles.}
  The W and Z bosons are identified as cross-peak lognormal residuals
  by their mass placement at bits 17--18.
  A derivation of their specific symmetry breaking from the codec
  structure (as opposed to the Higgs mechanism) is not yet available.

\item \textbf{Quark mass hierarchy.}
  The six quarks span a wider mass range than the leptons.
  Their placement in the $(n, m)$ classification and their log-mass
  integer assignments require a dedicated analysis.

\item \textbf{Many-fermion extension.}
  The present paper treats single-fermion states.
  The full Standard Model requires tensor products of single-fermion
  wavefunctions.
  Whether the $\mathrm{SU}(3) \times \mathrm{SU}(2) \times \mathrm{U}(1)$
  gauge structure emerges from the tensor product codec is the primary
  open question for the programme.

\item \textbf{Exact integer mass conjecture.}
  The near-integer values in Table~\ref{tab:masses} suggest
  $\log_2(m/m_e) \in \mathbb{Z}$ as a zeroth-order prediction,
  corrected by compression gains of order 0.1--0.4 bits.
  Whether these corrections follow from the spectral complexity formula
  $C_s$ and the Solomonoff weight $2^{-C_s}$ is an open calculation.
\end{enumerate}

%------------------------------------------------------------
\section{Conclusion}
\label{sec:conclusion}

We have proved that the Frobenius norm of the boson matrix satisfies
$\|B\| = \sqrt{(n-1)/n}$ for any $n$-site fermion in equal
superposition, with no dependence on phases or system size.
The winding number of the phase pattern determines the spin of the
emergent boson.
The pair $(n, m)$ classifies the known boson spectrum:

\[
\begin{array}{ll}
(n=2,\; m=1): & \text{Photon} \\
(n=3,\; m=0): & \text{Higgs candidate (scalar)} \\
(n=3,\; m=1): & \text{Gluon (colour-symmetric, confined)} \\
(n=4,\; m=2): & \text{Graviton candidate (spin-2)}
\end{array}
\]

Fractional charge $1/n$ per colour site emerges exactly.
Colour confinement follows from the mixed symmetry of the three-site
boson matrix.
The particle mass hierarchy organises into a near-integer sequence in
$\log_2(m/m_e)$ with step size $\approx 4$ bits, consistent with the
minimum fermionic codec unit.

These results suggest that the Standard Model particle spectrum is the
codec's natural output when the fermion site count $n$ ranges from 2
to 4 and the winding number $m$ ranges over its allowed values.
No interaction terms, coupling constants, or gauge symmetries were
introduced by hand.

%------------------------------------------------------------
\begin{thebibliography}{99}

\bibitem{meskanen2001}
J.~Meskanen,
\textit{Axiomatic Foundations of the Informational Universe},
Paper~0, IAME Collaboration (2001).

\bibitem{meskanen2026iv}
J.~Meskanen,
\textit{Emergence of Spacetime Geometry from Bitstring Compression},
Paper~IV, IAME Collaboration (2026).

\bibitem{meskanen2026v}
J.~Meskanen,
\textit{Emergence of General Relativity},
Paper~V, IAME Collaboration (2026).

\bibitem{meskanen2026vi}
J.~Meskanen,
\textit{The Wavefunction as Compression},
Paper~VI, IAME Collaboration (2026).

\bibitem{meskanen2026vii}
J.~Meskanen,
\textit{Emergence of Quantum Mechanics: Fermions, Bosons and Pauli Exclusion},
Paper~VII, IAME Collaboration (2026).

\end{thebibliography}

\end{document}
